\section{Threat Model}
\label{sec:threat}

\subsection{Adversarial setting}

QuasarCert is designed against an adversary who controls a
strict-minority subset of validators and observes all network
traffic. Specifically:

\begin{itemize}[leftmargin=2em,nosep]
\item \textbf{Byzantine fraction}: $f < 1/3$ in count and stake.
      Quorum threshold $\tau = \lfloor 2n/3 \rfloor + 1$.
\item \textbf{Static vs adaptive corruption}: theorems are stated for
      static corruption; adaptive corruption is handled via
      programmable-RO or secure-erasure hybrids
      \cite{fischlin2005} with $O(n \cdot Q_H)$ multiplicative loss.
      Production deployment uses Go-side memory zeroing for all
      per-signature randomness, satisfying the secure-erasure
      precondition.
\item \textbf{Network adversary}: full Dolev-Yao on the message
      passing layer; the adversary can drop, reorder, replay, and
      forge messages. Replay is prevented by round-digest binding
      (Theorem~\ref{thm:replay}); message authentication is via the
      cryptographic legs of the cert itself.
\item \textbf{Quantum adversary} ($\mathcal{A}_Q$): polynomial-time on
      a quantum Turing machine; runs Shor's algorithm on classical
      discrete logarithms; has QROM access \cite{bellare1993}; subject
      to Grover speedup against unstructured search. Under
      $\mathcal{A}_Q$, the BLS and Z-Chain legs break in polynomial
      time; the PQ legs remain hard
      (Corollary~\ref{cor:diversification}).
\item \textbf{Side-channel adversary}: passive timing / cache /
      power-analysis observer of validator processes. All threshold
      primitives are constant-time in their secret-dependent paths
      (documented per-primitive in \texttt{CONSTANT-TIME-REVIEW.md}
      / \texttt{ct/} subdirs).
\end{itemize}

\subsection{Threat surface boundaries}

\paragraph{In scope.} The QuasarCert verifier is in scope.
Specifically:

\begin{itemize}[leftmargin=2em,nosep]
\item Cert wire format and verification predicate
      (Definition~\ref{def:verify}).
\item Profile selection and profile invariants
      (Section~\ref{sec:profiles}).
\item Replay protection via round-digest binding
      (Theorem~\ref{thm:replay}).
\item Cross-primitive composition soundness
      (Theorem~\ref{thm:composition}).
\item Domain-separation contexts (Table~\ref{tab:domsep}).
\item Per-leg primitive unforgeability (delegated to per-primitive
      papers; see \cite{pulsarrepo,coronarepo,magnetarrepo,
      quasarcertsoundness}).
\end{itemize}

\paragraph{Out of scope (delegated to other layers).}

\begin{itemize}[leftmargin=2em,nosep]
\item Network-level partition tolerance: Lux uses a partial-synchrony
      assumption; full asynchronous safety / liveness is a property
      of the underlying Snow consensus family
      \cite{lp020,garaybackbone}, not the cert verifier.
\item Validator-set rotation: governed by the LSS framework
      \cite{luxlssmpc} and the per-primitive resharing protocols
      (\cite{pulsarrepo} for Pulsar, etc.). The cert sees only the
      post-resharing public keys.
\item Trusted setup for Z-Chain Groth16: governed by the
      Bowe-Gabizon-Miers ceremony \cite{bowegabizonmiers}; the
      verifying key is rooted on Z-Chain as
      $\mathsf{zchain\_vk\_root}$.
\item DKG-time security: each primitive runs its own DKG
      independently; the umbrella spec does not coordinate DKGs.
      Pulsar uses Pedersen-style DKG over $R_q$; Corona uses
      a similar lattice DKG; Magnetar uses byte-wise Shamir VSS
      over $\mathrm{GF}(257)$.
\item EVM precompile bindings: each PQ verifier is exposed via an
      EVM precompile (e.g., P3Q at slot \texttt{0x012205});
      precompile-layer profile gating uses
      \texttt{contract.RefuseUnderStrictPQ} in
      \texttt{luxfi/evm/precompile/}; cert verification is
      orthogonal to precompile verification.
\end{itemize}

\subsection{Failure modes}

We enumerate the failure modes and their consequences.

\paragraph{Single-primitive break.} An adversary breaks one PQ leg's
underlying assumption (e.g., a cryptanalytic break against M-LWE
specifically). On the Aurora profile, this would compromise the
Pulsar leg. The cert remains unforgeable because the Corona leg
(R-LWE) is independent and still valid (Corollary~\ref{cor:polaris}).
The compromised leg's primitive must be deprecated and replaced;
the wire format and other legs are unaffected.

\paragraph{Cross-lattice break.} An adversary breaks both M-LWE and
R-LWE (some breakthrough sieving / NTT-aware attack). On the
Aurora profile, this would compromise both Pulsar and Corona; the
profile no longer provides PQ security beyond the failing
$\mathrm{ZK}$ leg's MLDSA witness (which is also lattice-based).
On the Polaris profile, the Magnetar / hash-based leg remains
valid. This is the cross-family insurance that motivates Polaris.

\paragraph{Hash break.} An adversary breaks SHA-3 / SHAKE preimage
or target-collision resistance. On any profile, the round-digest
construction would be compromised. This is a catastrophic event
that affects all of cryptography, not just QuasarCert; the
remediation is to migrate to a different hash function and re-issue
all validator keys. The QuasarCert wire format makes this swap
non-disruptive in principle (replace $H_{\mathsf{round}}$ with a
new hash) but the per-primitive keys would need full rotation.

\paragraph{Trusted-setup compromise.} If an adversary compromises
the Z-Chain Groth16 trusted setup (e.g., the $\geq 1024$-contributor
ceremony \cite{bowegabizonmiers}), they can forge
$\pi^{\mathrm{ZK}}$. By Theorem~\ref{thm:composition}, the cert is
not forgeable because the other legs (BLS / Pulsar / Corona /
Magnetar) remain unbroken. The Z-Chain $\mathit{vk}_Z$ would be
rotated to a fresh ceremony output; the cert format is unchanged.

\paragraph{Aggregator-as-TCB (Pulsar v0.1 / Magnetar v0.1).} A
malicious aggregator can refuse to combine signatures (liveness
fault) but cannot produce a forgery (the aggregator's material is
purely derived from honest contributions; forging requires forging
at the underlying primitive). The fault is identifiable abort
\cite{cggmp21}; the consensus engine excludes the faulty aggregator
from future rounds.

\paragraph{Byzantine quorum.} An adversary controls $\geq 2/3$ of
validators. The cert remains valid (because the Byzantine signers
sign correctly under the protocol) but the consensus engine's
safety property is violated: the adversary can finalize conflicting
blocks. This is outside QuasarCert's scope --- QuasarCert is a
finality certificate, not a Byzantine-quorum gate. Quorum
intersection at $2f + 1$ of $3f + 1$ is proved by the Lean
\texttt{lean/Consensus/BFT.lean}~\cite{quasarcertsoundness}.

\subsection{Migration risk}

The PQ-migration era is the period during which a quantum adversary
becomes capable of breaking classical primitives in polynomial time
but the legacy classical primitives are still deployed. QuasarCert
addresses this risk in two ways:

\begin{itemize}[leftmargin=2em,nosep]
\item The PQ leg is always populated (profile invariant 2). A
      chain configured with the Pulsar profile is PQ-safe from day
      one; no migration is required.
\item The classical BLS leg is supplementary, not load-bearing.
      Under quantum attack, the BLS leg becomes trivially forgeable
      ($\epsilon_{\mathrm{BLS}} = O(1)$); by
      Corollary~\ref{cor:diversification}, the cert's
      unforgeability is bounded by the PQ leg's hardness, which
      remains intact.
\end{itemize}

The decision to keep BLS in the cert (rather than removing it
entirely; see \cite{pqfinalitynobls} for the BLS-removal analysis)
is a deployment-pragmatic one: BLS aggregate verify is
sub-millisecond, much faster than any PQ verifier; on classical
hardware, the cert verifies fast because BLS is checked first; on
quantum hardware, the cert is still safe because the PQ leg is
also verified.

The reverse decision --- remove BLS, run only PQ legs --- is also
implemented in \texttt{luxfi/consensus} as the ``pure PQ'' mode
(\texttt{pq-finality-no-bls}), gated behind a chain-config flag.
This is the migration target for chains that desire to retire all
classical cryptography from the consensus path.
